\documentclass[12pt]{article}
\usepackage[english]{babel}
\usepackage{natbib}
\usepackage{url}
\usepackage[utf8x]{inputenc}
\usepackage{amsmath}
\usepackage{graphicx}
\graphicspath{{images/}}
\usepackage{parskip}
\usepackage{fancyhdr}
\usepackage{vmargin}
\setmarginsrb{3 cm}{2.5 cm}{3 cm}{2.5 cm}{1 cm}{1.5 cm}{1 cm}{1.5 cm}

\usepackage{eurosym}                %simbolo dell'euro
\usepackage{listings}
\usepackage[usenames,dvipsnames,svgnames,table]{xcolor}
% Imposta lo spazio nella list of listing in modo simile alla list of figures/tables
%\makeatletter
%\let\my@chapter\@chapter
%\renewcommand*{\@chapter}{%
%  \addtocontents{lol}{\protect\addvspace{10pt}}%
%  \my@chapter}
%\makeatother


\definecolor{codegreen}{rgb}{0,0.6,0}
\definecolor{codegray}{rgb}{0.5,0.5,0.5}
\definecolor{backcolor}{rgb}{0.98,0.98,0.98}

\renewcommand{\lstlistingname}{Codice}% Listing -> codice
\renewcommand{\lstlistlistingname}{Elenco dei frammenti di codice}% List of Listings -> Frammenti di codice

\lstdefinestyle{mystyle}{
    backgroundcolor=\color{backcolor},   
    commentstyle=\color{Peach}\ttfamily,
    keywordstyle=\color{RoyalBlue},
    numberstyle=\tiny\color{codegray},
    stringstyle=\color{SeaGreen}\ttfamily,
    basicstyle=\footnotesize\ttfamily,
    breakatwhitespace=false,         
    breaklines=true,                 
    captionpos=b,                    
    keepspaces=true,                 
    numbers=left,                    
    numbersep=5pt,                  
    showspaces=false,                
    showstringspaces=false,
    showtabs=false,                  
    tabsize=2,
    frame=trbl, % draw a frame at the top, right, left and bottom of the listing
	frameround=ftff, % angolo in basso a destro curvo
	framesep=4pt, % quarter circle size of the round corners,
	inputencoding=utf8,
    extendedchars=true,
    literate={á}{{\'a}}1 {à}{{\`a}}1 {é}{{\'e}}1 {è}{{\`e}}1 {ù}{{\`u}}1 {ò}{{\`o}}1 {ì}{{\`i}}1,
    belowskip=1em,
    aboveskip=1em,
}

 
\lstset{style=mystyle}

\lstdefinelanguage{JavaScript}
{
  % list of keywords
  morekeywords={ true, false, catch, function, break,	new, class, extends, var, require, switch, return, import, if, while, for, this, View, Text, StyleSheet},
  sensitive=false, % keywords are not case-sensitive
  morecomment=[l]{//}, % l is for line comment
  morecomment=[s]{/*}{*/}, % s is for start and end delimiter
  morestring=[b]' % defines that strings are enclosed in double quotes
}

\lstdefinelanguage{JSON}
{
  % list of keywords
  morekeywords={string, boolean, int, Array, Node, Asset, AssetDetail, Filter, FilterItem},
  sensitive=false, % keywords are not case-sensitive
  morecomment=[l]{//}, % l is for line comment
  morecomment=[s]{/*}{*/}, % s is for start and end delimiter
  morestring=[b]" % defines that strings are enclosed in double quotes
}

\lstdefinelanguage{URM}
{
	% list of keywords
	morekeywords={ S, J, T, Z, I},
	sensitive=false, % keywords are not case-sensitive
	morecomment=[l]{//}, % l is for line comment
	morecomment=[s]{/*}{*/}, % s is for start and end delimiter
	morestring=[b]' % defines that strings are enclosed in double quotes
}

\lstdefinelanguage{RDFA}{
	language=html,
	sensitive=true, 
	alsoletter={<>=-},
	ndkeywords={
		% General
		=,
		% HTML attributes
		charset=, id=, width=, height=, property=, about=, rel=, rev=, prefix=, vocab=, content=, datatype=
	},  
	morecomment=[s]{<!--}{-->},
	tag=[s]
}

%\tightlist per compatibilità con pandoc
\providecommand{\tightlist}{%
	\setlength{\itemsep}{0pt}\setlength{\parskip}{0pt}}


\usepackage[labelfont=bf]{caption}

\usepackage[protrusion=true,expansion=true]{microtype} % Better typography
\usepackage{graphicx} % Required for including pictures
\usepackage{wrapfig} % Allows in-line images


\usepackage{subfig}
\usepackage{hyperref}
\usepackage{placeins}
\usepackage{sourcecodepro}
\usepackage{hyperref}                   % collegamenti ipertestuali

\usepackage[colorinlistoftodos,prependcaption]{todonotes} %todo

\usepackage{amsmath}
\usepackage{mathtools}

\usepackage{float}
\usepackage{algorithm}
\usepackage{algpseudocode} % https://en.wikibooks.org/wiki/LaTeX/Algorithms#Typesetting_using_the_algorithmicx_package
\usepackage{amsmath} 
\usepackage{amssymb}  %$\mathbb{N}$ per il simbolo dei numeri naturali 

\usepackage{enumerate} % permette di personalizzare enumerate
\usepackage{enumitem}
\usepackage{xmpincl}	%Aggiunge metadati sulla licenza CC
\usepackage{xspace}

\usepackage{bussproofs}

\usepackage{listings}
\usepackage{color}

\definecolor{dkgreen}{rgb}{0,0.6,0}
\definecolor{gray}{rgb}{0.5,0.5,0.5}
\definecolor{mauve}{rgb}{0.58,0,0.82}

\lstset{frame=tb,
  language=Java,
  aboveskip=3mm,
  belowskip=3mm,
  showstringspaces=false,
  columns=flexible,
  basicstyle={\small\ttfamily},
  numbers=none,
  numberstyle=\tiny\color{gray},
  keywordstyle=\color{blue},
  commentstyle=\color{dkgreen},
  stringstyle=\color{mauve},
  breaklines=true,
  breakatwhitespace=true,
  tabsize=3
}


\makeatletter
\renewcommand\@biblabel[1]{\textbf{#1.}} % Change the square brackets for each bibliography item from '[1]' to '1.'
\renewcommand{\@listI}{\itemsep=0pt} % Reduce the space between items in the itemize and enumerate environments and the bibliography

\renewcommand{\maketitle}{ % Customize the title - do not edit title and author name here, see the TITLE block below
	\begin{flushright} % Right align
		{\LARGE\@title} % Increase the font size of the title
		
		\vspace{50pt} % Some vertical space between the title and author name
		
		{\large\@author} % Author name
		\\\@date % Date
		
		\vspace{100pt} % Some vertical space between the author block and abstract
	\end{flushright}
}

%% breakablealgorithm http://tex.stackexchange.com/questions/33866/algorithm-tag-and-page-break
\makeatletter
\newenvironment{breakablealgorithm}
{% \begin{breakablealgorithm}
	\begin{center}
		\refstepcounter{algorithm}% New algorithm
		\hrule height.8pt depth0pt \kern2pt% \@fs@pre for \@fs@ruled
		\renewcommand{\caption}[2][\relax]{% Make a new \caption
			{\raggedright\textbf{\ALG@name~\thealgorithm} ##2\par}%
			\ifx\relax##1\relax % #1 is \relax
			\addcontentsline{loa}{algorithm}{\protect\numberline{\thealgorithm}##2}%
			\else % #1 is not \relax
			\addcontentsline{loa}{algorithm}{\protect\numberline{\thealgorithm}##1}%
			\fi
			\kern2pt\hrule\kern2pt
		}
	}{% \end{breakablealgorithm}
	\kern2pt\hrule\relax% \@fs@post for \@fs@ruled
\end{center}
}
\makeatother

\makeatletter % trattino con punto sopra
\newcommand{\dotminus}{\mathbin{\text{\@dotminus}}}

\newcommand{\@dotminus}{%
	\ooalign{\hidewidth\raise1ex\hbox{.}\hidewidth\cr$\m@th-$\cr}%
}
\makeatother

\DeclarePairedDelimiter{\ceil}{\lceil}{\rceil}
\DeclarePairedDelimiter{\floor}{\lfloor}{\rfloor}


\let\oldtext\text
\renewcommand{\text}[1]{\oldtext{\normalfont\sffamily #1}}
\newcommand{\false}{\text{ false }}
\newcommand{\true}{\text{ true }}
\newcommand{\fn}{\text{fn }}
\newcommand{\iif}{\text{ if }}
\newcommand{\then}{\text{ then }}
\newcommand{\eelse}{\text{ else }}
\newcommand{\Nat}{\text{ Nat }}
\newcommand{\Bool}{\text{ Bool }}

\newcommand{\IF}[3]{\ensuremath{\text{if}\ #1\ \text{then}\ #2\ \text{else}\ #3}}

\newcommand{\vbar}{\:|\:}
\newcommand{\proj}{.\_}

\newcommand{\throw}[1]{\ensuremath{\text{throw} \: #1}}
\newcommand{\trycatch}[2]{\ensuremath{\text{try} \: #1 \: \text{catch} \: #2}}
\newcommand{\TExc}{\ensuremath{T_{exc}}}
\newcommand{\myrule}[1]{\textsc{(#1)}}

\newcommand{\class}[3]{\text{class} \: #1 \: \text{extends} \: #2 \: \{\: #3 \:\}}
\newcommand{\construct}[3]{\ensuremath{#1(#2)\{\: #3\:\}}}
\newcommand{\this}{\text{this}}
\newcommand{\super}{\text{super}}

\newcommand{\method}[4]{\ensuremath{#1 \: #2(#3) \: \{ \: \text{return} \: #4;\: \} }}
\newcommand{\new}[2]{\ensuremath{\text{new}\: #1(  #2  )}}

\newcommand{\stype}{\ensuremath{<:}}
\newcommand{\fields}[1]{fields(#1)}
\newcommand{\mbody}[1]{mbody(#1)}
\newcommand{\mtype}[1]{mtype(#1)}
\newcommand{\IN}[1]{\text{in} \: #1}


% Shorcut per prooftree
\newcommand{\LL}[1]{\LeftLabel{\myrule{#1}}}
\newcommand{\AC}[1]{\AxiomC{#1}}
\newenvironment{bprooftree}
{\leavevmode\hbox\bgroup}
{\DisplayProof\egroup}

\title{Esercizi}                                % Title
\author{AALP}                               % Author
\date{3 Dec 2016}                                         % Date

\makeatletter
\let\thetitle\@title
\let\theauthor\@author
\let\thedate\@date
\makeatother

\pagestyle{fancy}
\fancyhf{}
\rhead{\theauthor}
\lhead{\thetitle}
\cfoot{\thepage}

\begin{document}

%%%%%%%%%%%%%%%%%%%%%%%%%%%%%%%%%%%%%%%%%%%%%%%%%%%%%%%%%%%%%%%%%%%%%%%%%%%%%%%%%%%%%%%%%

\begin{titlepage}
    \centering
    \vspace*{0.5 cm}
    \includegraphics[scale = 0.75]{logo.png}\\[1.0 cm]  % University Logo
    \textsc{\LARGE Aspetti Avanzati dei Linguaggi \newline\newline di Programmazione}\\[2.0 cm]  % University Name
    %\textsc{\Large CSX-325}\\[0.5 cm]               % Course Code
    %\rule{\linewidth}{0.2 mm} \\[0.4 cm]
    { \huge \bfseries \thetitle}\\
    \rule{\linewidth}{0.2 mm} \\[1.5 cm]
    
    \begin{minipage}{0.4\textwidth}
        \begin{flushleft} \large
            \emph{Autori:}\\
            Campese Stefano\\
            Enrico Savoca\\
        \end{flushleft}
    \end{minipage}~
    
    
    
    
    
    
    
\end{titlepage}

%%%%%%%%%%%%%%%%%%%%%%%%%%%%%%%%%%%%%%%%%%%%%%%%%%%%%%%%%%%%%%%%%%%%%%%%%%%%%%%%%%%%%%%%%

\tableofcontents
\pagebreak

%%%%%%%%%%%%%%%%%%%%%%%%%%%%%%%%%%%%%%%%%%%%%%%%%%%%%%%%%%%%%%%%%%%%%%%%%%%%%%%%%%%%%%%%%
  

\section{Il mini linguaggio funzionale (note 2)}
\subsection*{Esercizio 1.1}
La relazione di riduzione data definisce una strategia efficiente per la valutazione del termine if-then-else, che permette cioè di valutare unicamente il ramo scelto dalla valutazione della guardia
booleana. Ridefinire la semantica operazionale del linguaggio in modo che adotti una strategia non efficiente per il costrutto if-then-else, valutando entambi i rami del costrutto condizionale.

\subsubsection*{Svolgimento}
Per ottenere una semantica meno efficiente, si cambiano gli assiomi (IF-TRUE) e (IF-FALSE) e si aggiungono le regole (THEN) ed (ELSE). (IF) si mantiene tale.

\paragraph{Assiomi}
\begin{prooftree}
	\AxiomC{}
	\LeftLabel{\textsc{(IF-TRUE)}}
	\UnaryInfC{$if \: TRUE \: then \: v_1 \: else \: v_2 \rightarrow \: v_1$}
\end{prooftree}

\begin{prooftree}
	\AxiomC{}
	\LeftLabel{\textsc{(IF-FALSE)}}
	\UnaryInfC{$if \: FALSE \: then \: v_1 \: else \: v_2 \rightarrow \: v_2$}
\end{prooftree}

\paragraph{Regole}
\begin{prooftree}
	\AxiomC{$M_2 \rightarrow \: M_2'$}
	\LeftLabel{\textsc{(THEN)}}
	\UnaryInfC{$if \: v_1 \: then \: M_2 \: else \: M_3 \rightarrow \: 
	if \: v_1 \: then \: M_2' \: else \: M_3$}
\end{prooftree}

\begin{prooftree}
	\AxiomC{$M_3 \rightarrow \: M_3'$}
	\LeftLabel{\textsc{(ELSE)}}
	\UnaryInfC{$if \: v_1 \: then \: v_2 \: else \: M_3 \rightarrow \: 
	if \: v_1 \: then \: v_2 \: else \: M_3'$} 
\end{prooftree} 

\subsection*{Esercizio 1.4}
La valutazione, cioè l’esecuzione del programma, è deterministica. Se $M \rightarrow$ M' e
$M \rightarrow$ M'' allora M'=M''.\\
Dimostrare la proposizione precedente per induzione sulla struttura del termine M.

\subsubsection*{Dimostrazione}
Si dimostra per induzione sulla derivazione.

Casi base:
\begin{itemize}
	\item M=true \\
	\indent	Se true $\rightarrow$ M' e true $\rightarrow$ M'' $\implies$ M' = M''\\
	\indent Vero perchè true è un valore finale;
	
 \item M=false \\
	\indent	Se false $\rightarrow$ M' e false $\rightarrow$ M'' $\implies$ M' = M''\\
	\indent Vero perchè false è un valore finale;

\item M=fn x.M \\
	\indent	Se fn x.M $\rightarrow$ M' e fn x.M $\rightarrow$ M'' $\implies$ M' = M''\\
	\indent Vero perchè fn x.M è un valore finale;
	\item M=n \\
	\indent	Se n $\rightarrow$ M' e n $\rightarrow$ M'' $\implies$ M' = M''\\
	\indent Vero perchè n è un valore finale.
\end{itemize}
 
Passo induttivo:
\begin{itemize}
	\item Tesi: $M=M_1+M_2 \:\rightarrow\: M', \:M_1+M_2 \:\rightarrow$ M'' $\implies$ M'=M''\\
	\indent Per ipotesi induttiva: 
	\begin{itemize}
		\item Se $M_1 \:\rightarrow \: M_1'$ e $M_1 \:\rightarrow \: M_1''$ allora $M_1'=M_1''$;
		\item Se $M_2 \:\rightarrow \: M_2'$ e $M_2 \:\rightarrow \: M_2''$ allora $M_2'=M_2''$.
	\end{itemize} 
	3 Casi possibili:
	\begin{enumerate}[label=\alph*)]
		\item Derivazione tramite l'uso della regola (SUM-LEFT): 
		\begin{prooftree}
					\AxiomC{$M_1\:\rightarrow\:M_1'$}
					\LeftLabel{\textsc{(SUM-LEFT)}}
					\UnaryInfC{$M_1+M_2 \:\rightarrow \: M_1'+M_2=M'$}
				\end{prooftree} 
				\indent In questo caso $M_1'=M_1''$ per ipotesi 
						induttiva e quindi anche $M'=M''$;\\
				\indent La derivazione risulta deterministica applicando questa regola.
		\item Derivazione tramite l'uso della regola (SUM-RIGHT): 
		\begin{prooftree}
					\AxiomC{$M_2\:\rightarrow\:M_2'$}
					\LeftLabel{\textsc{(SUM-RIGHT)}}
					\UnaryInfC{$v_1+M_2 \:\rightarrow \: v_1+M_2'=M'$}
				\end{prooftree} 
				\indent In questo caso $M_2'=M_2''$ per ipotesi 
						induttiva e quindi anche $M'=M''$;\\
				\indent La derivazione risulta deterministica applicando questa regola.
		\item Derivazione tramite l'uso della regola (SUM): 
		\begin{prooftree}
					\AxiomC{con n'=$n_1+n_2$}
					\LeftLabel{\textsc{(SUM)}}
					\UnaryInfC{$n_1+n_2 \:\rightarrow \: n'$}
				\end{prooftree} 
				\indent Il risultato della somma è univoco e di conseguenza la derivazione è deterministica. 
	\end{enumerate}
	In tutti e 3 i casi, i soli possibili, la derivazione è deterministica, quindi la tesi 
	è dimostrata.\\
	
	
	\item Tesi: $M=M_1-M_2 \:\rightarrow\: M', \:M_1-M_2 \:\rightarrow$ M'' $\implies$ M'=M''\\
	\indent Per ipotesi induttiva: 
	\begin{itemize}
		\item Se $M_1 \:\rightarrow \: M_1'$ e $M_1 \:\rightarrow \: M_1''$ allora $M_1'=M_1''$;
		\item Se $M_2 \:\rightarrow \: M_2'$ e $M_2 \:\rightarrow \: M_2''$ allora $M_2'=M_2''$.
	\end{itemize} 
	3 Casi possibili:
	\begin{enumerate}[label=\alph*)]
		\item Derivazione tramite l'uso della regola (MINUS-LEFT): 
		\begin{prooftree}
					\AxiomC{$M_1\:\rightarrow\:M_1'$}
					\LeftLabel{\textsc{(MINUS-LEFT)}}
					\UnaryInfC{$M_1-M_2 \:\rightarrow \: M_1'-M_2=M'$}
				\end{prooftree} 
				\indent In questo caso $M_1'=M_1''$ per ipotesi 
						induttiva e quindi anche $M'=M''$;\\
				\indent La derivazione risulta deterministica applicando questa regola.
		\item Derivazione tramite l'uso della regola (MINUS-RIGHT): 
		\begin{prooftree}
					\AxiomC{$M_2\:\rightarrow\:M_2'$}
					\LeftLabel{\textsc{(MINUS-RIGHT)}}
					\UnaryInfC{$v_1-M_2 \:\rightarrow \: v_1-M_2'=M'$}
				\end{prooftree} 
				\indent In questo caso $M_2'=M_2''$ per ipotesi 
						induttiva e quindi anche $M'=M''$;\\
				\indent La derivazione risulta deterministica applicando questa regola.
		\item Derivazione tramite l'uso della regola (MINUS): 
		\begin{prooftree}
					\AxiomC{con n'=$n_1-n_2$}
					\LeftLabel{\textsc{(MINUS)}}
					\UnaryInfC{$n_1-n_2 \:\rightarrow \: n'$}
				\end{prooftree} 
				\indent Il risultato della differenza è univoco e di conseguenza la derivazione è deterministica. 
	\end{enumerate}
	In tutti e 3 i casi, i soli possibili, la derivazione è deterministica, quindi la tesi 
	è dimostrata.\\
	
	
	\item Tesi: $M=if\:M_1\:then\:M_2\:else\:M_3\:\rightarrow\: M',\:if\:M_1\:then\:M_2\:else\:M_3\:\rightarrow\: M''\:\implies$ M'=M''\\
	\indent Per ipotesi induttiva: 
	\begin{itemize}
		\item Se $M_1 \:\rightarrow \: M_1'$ e $M_1 \:\rightarrow \: M_1''$ allora $M_1'=M_1''$.
	\end{itemize}  
	3 Casi possibili:
	\begin{enumerate}[label=\alph*)]
		\item Derivazione tramite l'uso della regola (IF): 
		\begin{prooftree}
					\AxiomC{$M_1\:\rightarrow\:M_1'$}
					\LeftLabel{\textsc{(IF)}}
					\UnaryInfC{$if\:M_1\:then\:M_2\:else\:M_3\:\rightarrow\:if\:M_1'\:then\:M_2\:else\:M_3=M'$}
				\end{prooftree} 
				\indent In questo caso $M_1'=M_1''$ per ipotesi 
						induttiva e quindi anche $M'=M''$;\\
				\indent La derivazione risulta deterministica applicando questa regola.
		\item Derivazione tramite l'uso dell'assioma (IF-TRUE): 
		\begin{prooftree}
					\AxiomC{}
					\LeftLabel{\textsc{(IF-TRUE)}}
					\UnaryInfC{$if\:true\:then\:M_2\:else\:M_3\:\rightarrow\:M_2=M'$}
				\end{prooftree} 
				\indent In questo caso $M_2=M'=M''$ per applicazione dell'assioma e 
				la derivazione risulta deterministica.  
		\item Derivazione tramite l'uso della regola (IF-FALSE): 
		\begin{prooftree}
					\AxiomC{}
					\LeftLabel{\textsc{(IF-FAlSE)}}
					\UnaryInfC{$if\:false\:then\:M_2\:else\:M_3\:\rightarrow\:M_3=M'$}
				\end{prooftree} 
				\indent In questo caso $M_3=M'=M''$ per applicazione dell'assioma e 
				la derivazione risulta deterministica.  
	\end{enumerate} 
	In tutti e 3 i casi, i soli possibili, la derivazione è deterministica, quindi la tesi 
	è dimostrata.\\	
	
	
		\item Tesi: $M=M_1\:M_2 \:\rightarrow\: M', \:M_1\:M_2 \:\rightarrow$ M'' $\implies$ M'=M''\\
	\indent Per ipotesi induttiva: 
	\begin{itemize}
		\item Se $M_1 \:\rightarrow \: M_1'$ e $M_1 \:\rightarrow \: M_1''$ allora $M_1'=M_1''$;
		\item Se $M_2 \:\rightarrow \: M_2'$ e $M_2 \:\rightarrow \: M_2''$ allora $M_2'=M_2''$.
	\end{itemize}  
	2 Casi possibili:
	\begin{enumerate}[label=\alph*)]
		\item Derivazione tramite l'uso della regola (APP1): 
		\begin{prooftree}
					\AxiomC{$M_1\:\rightarrow\:M_1'$}
					\LeftLabel{\textsc{(APP1)}}
					\UnaryInfC{$M_1\:M_2 \:\rightarrow \: M_1'\:M_2=M'$}
				\end{prooftree} 
				\indent In questo caso $M_1'=M_1''$ per ipotesi 
						induttiva e quindi anche $M'=M''$;\\
				\indent La derivazione risulta deterministica applicando questa regola.
		\item Derivazione tramite l'uso della regola (APP2): 
		\begin{prooftree}
					\AxiomC{$M_2\:\rightarrow\:M_2'$}
					\LeftLabel{\textsc{(APP2)}}
					\UnaryInfC{$v_1\:M_2 \:\rightarrow \: v_1\:M_2'=M'$}
				\end{prooftree} 
				\indent In questo caso $M_2'=M_2''$ per ipotesi 
						induttiva e quindi anche $M'=M''$;\\
				\indent La derivazione risulta deterministica applicando questa regola. 
	\end{enumerate} 
	In tutti e 2 i casi, i soli possibili, la derivazione è deterministica, quindi la tesi 
	è dimostrata.\\	
	
	
		\item Tesi: $M=fn\:x.M\:v\:\rightarrow \:M\{x:=v\}=M'$, $fn\:x.M\:v\:\rightarrow \:M\{x:=v\}=M''\implies$ M'=M''\\  
 	Derivazione tramite l'uso dell'assioma (BETA): 
		\begin{prooftree}
					\AxiomC{}
					\LeftLabel{\textsc{(BETA)}}
					\UnaryInfC{$fn\:x.M\:v\:\rightarrow \:M\{x:=v\}=M'$}
				\end{prooftree} 
				\indent Tramite l'applicazione dell'assioma si ottiene M{x:=v}=M'=M''.\\
				\indent La derivazione risulta quindi deterministica e la tesi risulta dimostrata.\\	
\end{itemize}

La derivazione è determinstica all'applicazione di tutte le regole e gli assiomi del linguaggio, di conseguenza è deterministica l'esecuzione del programma.\\ 
\begin{flushright}
$\Box$
\end{flushright}

\subsection*{Esercizio 1.5}
Descrivere la valutazione del termine ((fn x.3) (fn y.y)) ((fn z.if z then 1 else 0) (false)).
Modificare le regole di valutazione in modo tale che, mantenendo una strategia call-by-value, il termine precedente evolva in un termine stuck in meno passi di riduzione. Scrivere le regole di valutazione della strategia call-by-name e valutare il termine precedente secondo questa strategia.\\

\subsubsection*{Svolgimento}
Valutazione tramite strategia call-by-value: \\
\textbf{passo 1}
\begin{prooftree} 
	\AxiomC{}
	\LeftLabel{\textsc{(BETA)}}
	\UnaryInfC{$(\fn x.3)(\fn y.y)\: \rightarrow \:3$}
	\LeftLabel{\textsc{(APP1)}}
	\UnaryInfC{$((\fn x.3)(\fn y.y))\:((\fn z.if\:z\:then\:1\:else\:0)(false))\: \rightarrow 
	\:3\:((\fn z.if\:z\:then\:1\:else\:0)(false))$}
\end{prooftree}
\textbf{passo 2}
\begin{prooftree} 
	\AxiomC{}
	\LeftLabel{\textsc{(BETA)}}
	\UnaryInfC{$(\fn z.if\:z\:then\:1\:else\:0)(false)\: \rightarrow \:if\:false\:then\:1\:else\:0$}
	\LeftLabel{\textsc{(APP2)}}
	\UnaryInfC{$3\:((\fn z.if\:z\:then\:1\:else\:0)(false))\: \rightarrow 
	\:3\:(if\:false\:then\:1\:else\:0)$}
\end{prooftree}
\textbf{passo 3}
\begin{prooftree} 
	\AxiomC{}
	\LeftLabel{\textsc{(IF-FALSE)}}
	\UnaryInfC{$if\:false\:then\:1\:else\:0\: \rightarrow \:0$}
	\LeftLabel{\textsc{(APP2)}}
	\UnaryInfC{$3\:(if\:false\:then\:1\:else\:0)\: \rightarrow 
	\:3\:0$}
\end{prooftree}

3 0 costituisce uno stuck: non esistono regole o assiomi per continuare l'esecuzione di tale programma. Esecuzione terminata dopo 3 passi di derivazione. \\
Si può mantenere la strategia CBV e terminare l'esecuzione in un numero minore di passi modificando (APP2):
 
\begin{prooftree}
	\AxiomC{$N\:\rightarrow\:N'$}
	\LeftLabel{\textsc{(APP2)}}
	\UnaryInfC{$(\fn x.M)\: N\:\rightarrow\:\fn x.M\:N'$} 
\end{prooftree}

In questo modo l'esecuzione del programma terminerebbe dopo un solo passo. Infatti, ci si ritroverebbe in una forma in cui non è possibile applicare alcuna regola o assioma (3 non è una funzione).\\

Con una strategia CBN (Call-by-name) non ho (APP2) e la funzione (BETA) è la seguente:

\begin{prooftree}
	\AxiomC{$\checkmark$}
	\LeftLabel{\textsc{(BETA)}} 
	\UnaryInfC{$(\fn x.M)\: N\: \rightarrow \:M\{x:=N\}$}
\end{prooftree}

Con questa regola mostriamo l'esecuzione del programma:\\

\textbf{passo 1}
\begin{prooftree} 
	\AxiomC{}
	\LeftLabel{\textsc{(BETA)}}
	\UnaryInfC{$(\fn x.3)(\fn y.y)\: \rightarrow \:3$}
	\LeftLabel{\textsc{(APP1)}}
	\UnaryInfC{$((\fn x.3)(\fn y.y))\:((\fn z.if\:z\:then\:1\:else\:0)(false))\: \rightarrow 
	\:3\:((\fn z.if\:z\:then\:1\:else\:0)(false))$}
\end{prooftree}

Stuck! Non esistono regole o assiomi applicabili: APP1 non è applicabile in quanto 3 non è funzione.



\subsection*{Esercizio 1.6}
Consideriamo le seguenti definizioni in Scala:\\
\begin{lstlisting}
def square(x:Int):Int = x*x  
def sumOfSquare(x:Int,y:Int):Int = square(x)+square(y)
\end{lstlisting}
Descrivere i passi di riduzione dell'espressione \textbf{sumOfSquare(3,4)} secondo una strategia call-by-value analoga a quella definita nella sezione precedente.\\ 
Descrivere inoltre la riduzione della stesa espressione secondo la strategia call-by-name. \\
Descrivere i passi di riduzione dell'espressione \textbf{sumOfSquare(3,2+2)} secondo le strategie call-by-value e call-by-name.

\subsubsection*{Svolgimento}
Caso \textbf{sumOfSquare(3,4)}.\\
CBV si comporta come CBN.\\

\textbf{passo 1}
\begin{prooftree} 
	\AxiomC{}
	\LeftLabel{\textsc{(BETA)}}
	\UnaryInfC{sumOfSquare(3, 4) $\rightarrow$ square(3) + square(4)} 
\end{prooftree}

\textbf{passo 2}
\begin{prooftree} 
	\AxiomC{}
	\LeftLabel{\textsc{(BETA)}}
	\UnaryInfC{square(3) $\rightarrow$ 3*3} 
	\LeftLabel{\textsc{(SUM-L)}}
	\UnaryInfC{square(3) + square(4) $\rightarrow$ 3*3 + square(4)} 
\end{prooftree}

\textbf{passo 3}
\begin{prooftree} 
	\AxiomC{}
	\LeftLabel{\textsc{(MULT)}}
	\UnaryInfC{3*3 $\rightarrow$ 9} 
	\LeftLabel{\textsc{(SUM-L)}}
	\UnaryInfC{3*3 + square(4) $\rightarrow$ 9 + square(4)} 
\end{prooftree}

\textbf{passo 4}
\begin{prooftree} 
	\AxiomC{}
	\LeftLabel{\textsc{(BETA)}}
	\UnaryInfC{square(4) $\rightarrow$ 4*4} 
	\LeftLabel{\textsc{(SUM-R)}}
	\UnaryInfC{9 + square(4) $\rightarrow$ 9 + 4*4} 
\end{prooftree}

\textbf{passo 5}
\begin{prooftree} 
	\AxiomC{}
	\LeftLabel{\textsc{(MULT)}}
	\UnaryInfC{4*4 $\rightarrow$ 16} 
	\LeftLabel{\textsc{(SUM-R)}}
	\UnaryInfC{9 + 4*4 $\rightarrow$ 9 + 16} 
\end{prooftree}

\textbf{passo 6}
\begin{prooftree} 
	\AxiomC{} 
	\LeftLabel{\textsc{(SUM)}}
	\UnaryInfC{9 + 16 $\rightarrow$ 25} 
\end{prooftree}

Caso \textbf{sumOfSquare(3,2+2)}\\
CBV:\\
\textbf{passo 1}
\begin{prooftree} 
	\AxiomC{}
	\LeftLabel{\textsc{(SUM)}}
	\UnaryInfC{2+2 $\rightarrow$ 4} 
	\LeftLabel{\textsc{(APP2)}}
	\UnaryInfC{sumOfSquare(3, 2+2) $\rightarrow$ sumOfSquare(3, 4)} 
\end{prooftree}

I passi successivi sono gli stessi del caso precedente, quindi i passi di esecuzione sono 7.\\

CBN:\\
\textbf{passo 1}
\begin{prooftree} 
	\AxiomC{} 
	\LeftLabel{\textsc{(BETA)}}
	\UnaryInfC{sumOfSquare(3, 2+2) $\rightarrow$ square(3) + square(2+2)} 
\end{prooftree}

\textbf{passo 2}
\begin{prooftree} 
	\AxiomC{} 
	\LeftLabel{\textsc{(BETA)}}
	\UnaryInfC{square(3) $\rightarrow$ 3*3}
	\LeftLabel{\textsc{(SUM-L)}}
	\UnaryInfC{square(3) + square(2+2) $\rightarrow$ 3*3 + square(2+2)} 
\end{prooftree}

\textbf{passo 3}
\begin{prooftree} 
	\AxiomC{} 
	\LeftLabel{\textsc{(MULT)}}
	\UnaryInfC{3*3 $\rightarrow$ 9}
	\LeftLabel{\textsc{(SUM-L)}}
	\UnaryInfC{3*3 + square(2+2) $\rightarrow$ 9 + square(2+2)} 
\end{prooftree}

\textbf{passo 4}
\begin{prooftree} 
	\AxiomC{} 
	\LeftLabel{\textsc{(BETA)}}
	\UnaryInfC{square(2+2) $\rightarrow$ (2+2)*(2+2)}
	\LeftLabel{\textsc{(SUM-R)}}
	\UnaryInfC{9 + square(2+2) $\rightarrow$ 9 + (2+2)*(2+2)} 
\end{prooftree}

\textbf{passo 5}
\begin{prooftree} 
	\AxiomC{} 
	\LeftLabel{\textsc{(SUM)}}
	\UnaryInfC{2+2 $\rightarrow$ 4}
	\LeftLabel{\textsc{(MULT-L)}}
	\UnaryInfC{(2+2)*(2+2) $\rightarrow$ (4)*(2+2)}
	\LeftLabel{\textsc{(SUM-R)}}
	\UnaryInfC{9 + (2+2)*(2+2) $\rightarrow$ 9 + (4)*(2+2)} 
\end{prooftree}

\textbf{passo 6}
\begin{prooftree} 
	\AxiomC{} 
	\LeftLabel{\textsc{(SUM)}}
	\UnaryInfC{2+2 $\rightarrow$ 4}
	\LeftLabel{\textsc{(MULT-R)}}
	\UnaryInfC{4*(2+2) $\rightarrow$ 4*4}
	\LeftLabel{\textsc{(SUM-R)}}
	\UnaryInfC{9 + 4*(2+2) $\rightarrow$ 9 + 4*4} 
\end{prooftree}

\textbf{passo 7}
\begin{prooftree} 
	\AxiomC{}  
	\LeftLabel{\textsc{(MULT)}}
	\UnaryInfC{4*4 $\rightarrow$ 16}
	\LeftLabel{\textsc{(SUM-R)}}
	\UnaryInfC{9 + 4*4 $\rightarrow$ 9 + 16} 
\end{prooftree}

\textbf{passo 8}
\begin{prooftree} 
	\AxiomC{}   
	\LeftLabel{\textsc{(SUM)}}
	\UnaryInfC{9 + 16 $\rightarrow$ 25} 
\end{prooftree}

I passi di esecuzione sono 8 in questo caso (CBN).

\subsection*{Esercizio 1.7}
Si consideri la seguente definizione in Scala:

\begin{lstlisting}
def test(x:Int, y:Int):Int = x*x
\end{lstlisting}
confrontare la velocit\`a (cio\`e il numero di passi) di riduzione delle seguenti espressioni secondo le strategie CBV e CBN, indicando quale delle due \`e pi\`u veloce. 
\begin{enumerate}
	\item test(2,3)
	\item test(3+4,8)
	\item test(7,2*4)
	\item test(3+4,2*4)
\end{enumerate}

\subsubsection*{Svolgimento}
\begin{enumerate}
	\item CBV=CBN: entrambi terminano l'esecuzione in 2 passi; 
	\item CBV<CBN: CBV termina prima perch\`e valuta l'espressione 3+4 subito (una sola volta). CBN invece invoca prima la funzione, sostituendo ad x l'espressione. Essendo x presente 2 volte nella funzione, la valutazione viene effettuata 2 volte. 
	\item CBV>CBN: CBN termina prima perch\`e non valuta l'espressione 2*4 che non \`e utilizzata nella funzione test.
	\item CBV=CBN: entrambi terminano l'esecuzione in 4 passi; la differenza sta nel fatto che con CBV si svolgono prima le due operazioni (somma e moltiplicazione) e poi si chiama la funzione. Con CBN invece, si evita di calcolare 2*4, tuttavia si calcola 3+4 due volte.
\end{enumerate}



\subsection*{Esercizio 1.8}
Definire in Scala una funzione and che si comporta come il costrutto logico $x\&\&y$. 

\subsubsection*{Svolgimento}

\begin{lstlisting}
// and logico. b1 e b2 passati by-name 
def and(b1:=> Boolean, b2:=> Boolean): Boolean =
	if(!b1) false
	else if(b2) true
		 else false; 
\end{lstlisting}

Il parametro b1 viene passato by-name nella nostra soluzione. Tuttavia, esso pu\`o essere passato anche by-value visto che in ogni caso viene valutato. 











\section{I tipi semplici (note 3)}
\subsection*{Esempi di derivazione svolti}
 
\begin{enumerate}[label=\alph*)]
	\item $\emptyset\:\vdash$  if true then 5 + 7 else 2 :Nat
	\item $\emptyset\:\vdash$  $\fn$ x:T.x :T$\rightarrow$T
	\item $\emptyset\:\vdash$  ($\fn$ x:Bool.x) true :Bool
	\item f:Bool$\rightarrow$ Bool $\vdash$ f (if false then true else false) :Bool
	\item f:Bool$\rightarrow$ Bool $\vdash\:\fn$ x:Bool.f(if x then false else x) :Bool$\rightarrow$Bool
\end{enumerate}

\subsubsection*{Svolgimento}

\begin{enumerate}[label=\alph*), leftmargin=*]
	\item $\emptyset\:\vdash$  if true then 5 + 7 else 2 :Nat\\
		\scalebox{.75}{
		\parbox{1cm}{
    	\begin{prooftree}
			\AxiomC{$\checkmark$}
			\LeftLabel{\textsc{(TRUE)}}
			\UnaryInfC{$\emptyset\:\vdash$ true: Bool}
			\AxiomC{$\checkmark$}
			\LeftLabel{\textsc{(NAT)}}
			\UnaryInfC{$\emptyset\:\vdash$ 5 :Nat}
			\AxiomC{$\checkmark$}
			\LeftLabel{\textsc{(NAT)}}
			\UnaryInfC{$\emptyset\:\vdash$ 7 :Nat}
			\LeftLabel{\textsc{(SUM)}}
			\BinaryInfC{$\emptyset\:\vdash$ 5 + 7 :Nat}
			\AxiomC{$\checkmark$}
			\LeftLabel{\textsc{(NAT)}}
			\UnaryInfC{$\emptyset\:\vdash$  2 :Nat}
			\LeftLabel{\textsc{(IF-THEN-ELSE)}}
			\TrinaryInfC{$\emptyset\:\vdash$  if true then 5 + 7 else 2 :Nat}
		\end{prooftree}}}

    \vspace*{1 cm}

	\item $\emptyset\:\vdash$  $\fn$ x:T.x :T$\rightarrow$T  \\
		\scalebox{.75}{
		\parbox{1cm}{
		\begin{prooftree} 
			\AxiomC{x:T $\in$ x:T}
			\LeftLabel{\textsc{(VAR)}}
			\UnaryInfC{x:T $\vdash$  x :T}
			\LeftLabel{\textsc{(FUN)}}
			\UnaryInfC{$\emptyset\:\vdash$  $\fn$ x:T.x :T$\rightarrow$T}
		\end{prooftree}}
		}
	 
	 
    \vspace*{1 cm}
	\item $\emptyset\:\vdash$  ($\fn$ x:Bool.x) true :Bool\\
		\scalebox{.75}{
		\parbox{1cm}{
		\begin{prooftree} 
			\AxiomC{x:Bool $\in$ x:Bool}
			\LeftLabel{\textsc{(VAR)}}
			\UnaryInfC{x:Bool $\vdash$  x :Bool}
			\LeftLabel{\textsc{(FUN)}}
		    \UnaryInfC{$\emptyset\:\vdash\:\fn$ x:Bool.x :Bool$\rightarrow					$Bool}
			\AxiomC{$\checkmark$}
			\LeftLabel{\textsc{(TRUE)}}
			\UnaryInfC{$\emptyset\:\vdash$ true :Bool}
			\LeftLabel{\textsc{(APP)}}
			\BinaryInfC{$\emptyset\:\vdash$  ($\fn$ x:Bool.x) true :Bool}
		\end{prooftree}}
		}
		
	
    \vspace*{1 cm}	
	\item f:Bool$\rightarrow$ Bool $\vdash$ f (if false then true else false) :Bool\\
		Sia $\Gamma$ = f:Bool$\rightarrow$Bool \\
		\scalebox{.75}{
		\parbox{1cm}{
		\begin{prooftree}  
			\AxiomC{f:Bool$\rightarrow$ Bool $\in$ $\Gamma$}
			\LeftLabel{\textsc{(VAR)}}
			\UnaryInfC{$\Gamma$ $\vdash$ f :Bool $\rightarrow$ Bool} 
			\AxiomC{$\checkmark$}
			\LeftLabel{\textsc{(FALSE)}}
			\UnaryInfC{$\Gamma$ $\vdash$ false :Bool}
			\AxiomC{$\checkmark$}
			\LeftLabel{\textsc{(TRUE)}}
			\UnaryInfC{$\Gamma$ $\vdash$ true :Bool}
			\AxiomC{$\checkmark$}
			\LeftLabel{\textsc{(FALSE)}}
			\UnaryInfC{$\Gamma$ $\vdash$ false :Bool}
			\LeftLabel{\textsc{(IF-THEN-ELSE)}}
			\TrinaryInfC{$\Gamma$ $\vdash$ if false then true else false :Bool}
			\LeftLabel{\textsc{(APP)}}
			\BinaryInfC{$\Gamma$ $\vdash$ f (if false then true else 							false):Bool}
		\end{prooftree}}
		} 
	
	
    \vspace*{1 cm}
	\item f:Bool$\rightarrow$ Bool $\vdash\:\fn$ x:Bool.f(if x then false else x) :Bool$\rightarrow$Bool 
		Sia $\Gamma$ = f:Bool$\rightarrow$Bool \\
		\scalebox{.75}{
		\parbox{1cm}{
		\begin{prooftree}  
			\AxiomC{f:Bool $\rightarrow$ Bool $\in$ $\Gamma$, x:Bool}
			\LeftLabel{\textsc{(VAR)}}
			\UnaryInfC{$\Gamma$, x:Bool $\vdash$ f: Bool $\rightarrow$ Bool} 
			\LeftLabel{\textsc{(FUN)}}
			\UnaryInfC{$\Gamma$ $\vdash\:\fn$ x:Bool.f: ? $\rightarrow$ Bool 
			$\rightarrow$ Bool} 
			\AxiomC{x:Bool $\in$ $\Gamma$}
			\LeftLabel{\textsc{(VAR)}}
			\UnaryInfC{$\Gamma$ $\vdash$ x :?}
			\AxiomC{$\checkmark$}
			\LeftLabel{\textsc{(TRUE)}}
			\UnaryInfC{$\Gamma$ $\vdash$ true :Bool}
			\AxiomC{$\checkmark$}
			\LeftLabel{\textsc{(FALSE)}}
			\UnaryInfC{$\Gamma$ $\vdash$ false :Bool}
			\LeftLabel{\textsc{(IF-THEN-ELSE)}}
			\TrinaryInfC{$\Gamma$ $\vdash$ (if x then false else x):? 
			$\rightarrow$ Bool}
			\LeftLabel{\textsc{(APP)}}
			\BinaryInfC{$\Gamma\:\vdash\:\fn$ x:Bool.f(if x then false 
			else x)
			 :Bool$\rightarrow$Bool}
		\end{prooftree}}
		} \\
		Quindi ?=Bool.	
\end{enumerate}
 
    \vspace*{1 cm}

\subsection*{Esercizio 2.1}
Trovare un contesto $\Gamma$ tale che $\Gamma\:\vdash$ f x y : Bool sia derivabile.

\subsubsection*{Svolgimento}
 	\begin{prooftree} 
		\AxiomC{f:$T_2\:\rightarrow\:T_1\:\rightarrow$ Bool $\in\:\Gamma$}
		\LeftLabel{\textsc{(VAR)}}
		\UnaryInfC{$\Gamma\:\vdash$ f :$T_2\:\rightarrow\:T_1\:\rightarrow$
		 Bool}
		\AxiomC{x:$T_2\:\in$ x:$T_2$}
		\LeftLabel{\textsc{(VAR)}}
		\UnaryInfC{$\Gamma\:\vdash$ x :$T_2$}
		\LeftLabel{\textsc{(APP)}}
		\BinaryInfC{$\Gamma\:\vdash$ f x :$T_1\:\rightarrow$ Bool}
		\AxiomC{y :$T_1$ $\in$ $\Gamma$}
		\LeftLabel{\textsc{(VAR)}}
		\UnaryInfC{$\Gamma\:\vdash$ y :$T_1$}
		\LeftLabel{\textsc{(APP)}}
		\BinaryInfC{$\Gamma\:\vdash$ f x y :Bool}
	\end{prooftree} 
	
	Esiste un contesto per questo programma, in quanto il programma si pu\`o ottenere applicando le regole e costruendo l'albero di derivazione. \\
	Fanno parte del contesto le seguenti regole di tipo:
\begin{itemize}
	\item f:$T_2\:\rightarrow\:T_1\:\rightarrow$ Bool
	\item x :$T_2$
	\item y :$T_1$
\end{itemize}

\vspace*{2 cm}
Caso f (x y):\\
 	\begin{prooftree} 
		\AxiomC{f :$T_1\:\rightarrow$ Bool $\:\in\:\Gamma$}
		\LeftLabel{\textsc{(VAR)}}
		\UnaryInfC{$\Gamma\:\vdash$ f :$T_1\:\rightarrow$ Bool}
		\AxiomC{x:$T_2\:\rightarrow\:T_1\:\in\:\Gamma$}
		\LeftLabel{\textsc{(VAR)}}
		\UnaryInfC{$\Gamma\:\vdash$ x :$T_2\:\rightarrow\:T_1$}
		\AxiomC{y:$T_2\:\in\:\Gamma$}
		\LeftLabel{\textsc{(VAR)}}
		\UnaryInfC{$\Gamma\:\vdash$ y :$T_2$}
		\LeftLabel{\textsc{(APP)}}
		\BinaryInfC{$\Gamma\:\vdash$ x y :$T_1$}
		\LeftLabel{\textsc{(APP)}}
		\BinaryInfC{$\Gamma\:\vdash$ f (x y) :Bool}
	\end{prooftree} 
	
	Esiste un contesto anche per questo programma. \\
	Fanno parte del contesto le seguenti regole di tipo:
\begin{itemize}
	\item f:$T_1\:\rightarrow\:$ Bool
	\item x :$T_2\:\rightarrow\:T_1$
	\item y :$T_2$
\end{itemize}


\subsection*{Esercizio 2.2}
%Il giudizio $\Gamma\:\vdash$ x x : T \`e derivabile? Se s\`i, trovare una derivazione per qualche $\Gamma$, T, altrimenti provare che non \`e derivabile.

\subsubsection*{Svolgimento}
 	\begin{prooftree} 
		\AxiomC{$\Gamma\:\vdash$ x :$T_1\:\rightarrow\:$T}
		\AxiomC{$\Gamma\:\vdash$ x :$T_1$}
		\LeftLabel{\textsc{(APP)}}
		\BinaryInfC{$\Gamma\:\vdash$ x x :T}
	\end{prooftree} 

No, non \`e derivabile. Nel nostro sistema di tipi x non pu\`o assumere contemporaneamente il tipo T e il tipo $T_1\:\rightarrow\:$T. Risulterebbe derivabile se fosse presente il tipo ricorsivo nel sistema di tipi. \\



\section{I Tipi Semplici (note 45)}
 
\subsubsection*{Esercizio 3.2}
Provare che ogni sottotermine di un termine ben tipato e' ben tipato.
\subparagraph{Svolgimento}
Per la risoluzione di questo esericizio e' necessario utilizzare il metodo dell'induzione sui passi di derivazione, ovvero sull'altezza dell'albero di derivazione.\\
Se $M$ \`e un termine ben tipato vuol dire che il giudizio di tipo $\Gamma$ $\vdash$ $M$ : $T$ \`e derivabile con un albero di derivazione di altezza $h$ e quindi l'induzione viene fatta su $h$.


\begin{description}

 \item[Caso Base $h=1$] A questo livello il termine e il sottotermine coincidono. Questo perch\`e l'albero consiste nell'applicazione di un assioma e quindi il termine $M$ non contiene sotto-termini, quindi la propriet\`a \`e verificata.

 \item[Caso Induttivo $h+1$ ] Questo caso \`e composto dall'applicazione della regola e dai vari sotto-alberi che provano le ipotesi della regola.\\I sotto-alberi in questione, riguardano il giudizio dei sotto-termini di $M$ e sono di altezza inferiore rispetto a questo albero ($h$). Possiamo quindi applicare l'ipotesi induttiva per definire ch ei sotto-termini dei sotto-termini di $M$ sono ben tipati e quindi di conseguenza lo sono anche i sotto-termini di $M$. A questo punto la propriet\`a risulta essere verificata.
 
 \end{description}
 
 \`E possibile inoltre, utilizzare un altra dimostrazione che sfrutta il \emph{la struttura di $M$} e il \emph{Lemma di Inversione}.
 
 
 \begin{itemize}[label=$\star$]

 \item $M$ $\equiv$ $v$:  In questo caso $M$ \`e un valore, come ad esempio $true$, $false$ o $n$ e quindi, poich\`e questi sono termini privi di sotto-termini la propriet\`a risulta essere verificata.

 \item $M$ $\equiv$ $A+B$:  In questo caso, per il Lemma di Inversione se $M$ \`e ben tipato, allora lo sono anche $A$ e$B$ ( tipo $Nat$ )
 
 \item $M$ $\equiv$ $A-B$:  questo caso \`e analogo al precendente

 \item $M$ $\equiv$ $if$ $C$ $then$ $A$ $else$ $B$:  In questo caso, per il Lemma di Inversione se $M$ \`e ben tipato, allora $C$ \`e ben tipato con tipo $Bool$ mentre $A$ e $B$ sono ben tipati con tipo $T$
 
  \item $M$ $\equiv$ $A$ $B$: In questo caso, per il Lemma di Inversione se $M$ \`e ben tipato, allora $A$ \`e ben tipato con tipo $T_1 \rightarrow T$ mentre $B$ \`e ben tipato con tipo $T$
  
    \item $M$ $\equiv$ $fn$ $x.N$: In questo caso, per il Lemma di Inversione $N$ \`e ben tipato con tipo $T$
 \end{itemize}
 
 
\subsubsection*{Esercizio 3.12} 
Dimostrare subject reduction per induzione sulla derivazione di $\Gamma\:\vdash$ M : T.

\subparagraph*{Svolgimento}

Si vuole quindi dimostrare che:

\begin{center}
	Se $\Gamma \vdash{} M : T$ e $M \to{} M'$, allora $\Gamma \vdash{} M' : T$.
\end{center}
%La dimostrazione per induzione va fatta sul numero di passi che $M$ portano a $M'$.
%Dobbiamo ragionare su quella che e' la lunghezza dei cammini della derivazione:

%Caso Base l = 0:\\
%Otteniamo che $M$ = $M'$ e la tesi e' trivialmente vera.

%Caso Induttivo l+1:\\
%In questo passo ho che $M \rightarrow^{*}_l M_1 \rightarrow M'$, e per ipotesi induttiva ho che $\Gamma \vdash M_1 : T$.\\
%A questo punto ho che $M_1 \rightarrow M'$ e quindi posso applicare subject-reduction per ottenere che $\Gamma \vdash M' : T$.

%Ovviamente grazie ai teoremi di Sostituzione, Preservazione e Progressione so per certo che qualsiasi sia il contesto %il risultato non cambia

\begin{description}


 \item[True] $\Gamma \vdash{} M : T$ \'e $\Gamma \vdash{} \true{} :
  \Bool$.
  Il teorema \'e vacuamente vero per questa derivazione siccome
  $\not \exists{M'}. M \to{} M'$.
  
 \item[False]
  $\Gamma \vdash{} M : T$ \'e $\Gamma \vdash{} \false{} : \Bool$.
  Analogo a True.

\item[Nat]
  $\Gamma \vdash{} M : T$ \'e $\Gamma \vdash{} n : \Nat$. Analogo a True.

\item[Var] $\Gamma \vdash{} M : T$ \'e $\Gamma \vdash{} x : T$.
  Analogo a True.

\item[Fun] $\Gamma \vdash{} M : T$ \'e
  $\Gamma \vdash{} fn x\:T_1 C : T_1 \to{} T_2$. Analogo a True.
  
  \item[Sum]
  $\Gamma \vdash{} M : T$ \'e $\Gamma \vdash{} A + B : \Nat$.
 
  Per quanto assunto dalla regola di tipo \emph{Sum}, sappiamo per certo che
  anche $A$ e $B$ hanno tipo $\Nat$ sotto contesto $\Gamma$.

  Supponiamo che $A + B \to{} M'$ (sappiamo che $M'$ \'e unico poich\'e l'esecuzione
   \'e  deterministica).Tale riduzione pu\'o essere:

  \begin{itemize}
    \item \emph{Sum}: $M' \equiv{} n$. Per la regola di tipo
      \emph{Nat}, si ha che $T \equiv{} \Nat$ e che
      $\Gamma{} \vdash{} M' : \Nat$;
    \item \emph{Sum-Left}: $M' \equiv{} A' + B$. 
    Per gli assunti della regola $\emph{Sum}$ ho che $\Gamma \vdash A : \Nat$ e per $\emph{Sum-Left}$ $A \to A'$.
    Posso quindi applicare l'ipotesi induttiva\footnote{perch \'e  la derivazione $\Gamma \vdash A :\Nat$ richiede meno passaggi}
    per ottenere il giudizio $ \Gamma \vdash A' : \Nat$.
    Sono quindi soddisfatte le premesse della regola $\emph{Sum}$ per $M'$ e quindi  $\Gamma{} \vdash{} M' : \Nat \equiv T$;

    \item \emph{Sum-Right}: $M' \equiv{} v + B'$. Il ragionamento  \'e  analogo al caso $\emph{Sum-Left}$.
  \end{itemize}

\item[Minus]
  $\Gamma \vdash{} M : T$ \'e $\Gamma \vdash{} A - B : \Nat$. Analogo a
  \emph{Sum}.

\item[IfThenElse] $\Gamma \vdash{} M : T$ \'e 
  $\Gamma \vdash{} \IF{M_1}{M_2}{M_3}$.
  
  Dal giudizio di tipo abbiamo che:

  \begin{itemize}
  	\item $\Gamma \vdash M_1 : \Bool$;
    \item $\Gamma \vdash{} M_2 : T$;
    \item $\Gamma \vdash{} M_3 : T$.
  \end{itemize}

  Supponiamo che $\IF{M_1}{M_2}{M_3} \to{} M'$ (sappiamo che $M'$ \'e unico poich\'e l'esecuzione
  \'e deterministica). Tale riduzione pu\'o essere:

  \begin{itemize}
    \item \myrule{If}: $M' \equiv{} \Gamma \vdash{} \mbox{if } M'_1
      \mbox{ then } M_2 \mbox{ else } M_3$. Dato che il giudizio $\Gamma \vdash M_1 : \Bool$ richiede una derivazione in meno e che $M_1 \to M_1'$, possiamo applicare l'ipotesi induttiva per ottenere il giudizio $\Gamma \vdash M_1' : \Bool$. Si ha quindi che la regola di tipo \myrule{IfThenElse} continua a valere per il termine $M'$ e dato che i termini $M_2$ e $M_3$ non sono cambiati, anche $M'$ ha tipo $T$;
    \item \myrule{If-True}: $M' \equiv{} M_2$. $M_2$ ha tipo $T$ e quindi anche $M'$ ha tipo $T$;
    \item \myrule{If-False}: $M' \equiv{} M_3$.  Analogo a \myrule{If-True}
  \end{itemize}
  
  
\item[App] $\Gamma \vdash{} M : T$ \'e  $\Gamma \vdash{} A \: B : T_2$.
  Dal giudizio di tipo abbiamo che:

  \begin{itemize}
    \item $\Gamma \vdash{} A : T_1 \to{} T_2$;
    \item $\Gamma \vdash{} B : T_1$.
  \end{itemize}

  Supponiamo che $A \: B \to{} M'$ (sappiamo che $M'$ \'e  unico poich\'e  l'esecuzione
  \'e  deterministica). Tale riduzione pu\'o essere:

  \begin{itemize}
    \item \myrule{Beta}: $M' \equiv{} \Gamma \vdash{}
      Z[x \coloneqq{} B]$\footnote{dato $A \equiv{} FN{x}{T_1}Z$}.
      Per la regola di tipo $\myrule{Fun}$ ho che $Z : T_2$ e quindi, per
       il Substitution Lemma, $\Gamma \vdash{} Z[x \coloneqq{} B] : T_2$. Per
      questo motivo si ha che $\Gamma \vdash{} M' : T_2$ con $T_2 \equiv T$;
    \item \myrule{App-1}: $M' \equiv{} A'\: B$. Dato che il giudizio $\Gamma \vdash A : T_1 \to T_2$ richiede una derivazione con meno passaggi e che $A \to A'$ posso applicare l'ipotesi induttiva per ottenere il giudizio $A' : T_1 \to T_2$. Il termine $M'$ soddisfa le premesse della regola di tipo \myrule{App} e quindi $M'$ ha tipo $T_2 \equiv T$.

    \item \myrule{App-2}: $M' \equiv{} v\: B'$. Analogo ad \myrule{App-1}.
  \end{itemize}

 
\end{description}




\subsubsection*{Esercizio 3.13}

Vale l'opposto di subject reduction, i.e. se $\Gamma\:\vdash$ $M'\: :  $T e $M\:\rightarrow\:$M', allora $\Gamma\:\vdash$ $M$ : $T$ (detto subject expansion)? Dimostrarlo oppure dare un controesempio.
\subparagraph*{Svolgimento}
La risposta e' no:

E sufficiente pensare a $M \rightarrow M_1$ con $\Gamma\vdash$ $M_1$ : $\sigma$ e che 
$\Gamma\nvdash$ $M$: $\sigma$\\
Quanto riportato sopra pu\`o essere semplificato tramite questo esempio:\\
Supponimao di avere $M_1$ $\equiv$ $A$ : $T$ e $M$ $\equiv$ $if$ $true$ $then$ $A$ $else$ $B$, ovvero il caso in cui $M$ $\rightarrow$ $M_1$ per la regola (If-True).\\
In questo caso la derivazione di $\Gamma$ $\vdash$ $M$ : $T$ ri riesce ad ottenere usando la regaola $IfThenElse$, tuttavia questo impone che nel contesto ci siano questi giudizi:
\begin{itemize}
\item $\Gamma$ $\vdash$ $true$:$Bool$ che \`e un assioma
\item $\Gamma$ $\vdash$ $A$:$T$ soddisfatto per l'ipotesi 
\item  $\Gamma$ $\vdash$ $B$:$T$ putroppo questa informazione non si riesce a repire dall'ipotesi ne dal contesto e pertatnto non c'\`e garanzia sull'uguaglianza di tipi di $A$ e $B$\\
Pertanto possiamo creare un termine $N$ come $if$ $true$ $then$ $4$ $else$ $false$, da cui si riesce ad attenere $M'$ $\equiv$ $A$ $\equiv$ $4$. Questo termine nonostante esegua risulta essere non ben-tipato.


\end{itemize}







\subsubsection*{Esercizio 3.14}
Se al posto della regola $(APP)\:$ si definisce la regola seguente:

\begin{prooftree} 
	\AxiomC{$\Gamma\:\vdash$ $M\: :  $T$\:\rightarrow\:$T}
	\AxiomC{$\Gamma\:\vdash$ $N\ :$T}
	
	\LeftLabel{\textsc{(APP')}}
	\BinaryInfC{$\Gamma\:\vdash$ $M$ $N$ : $T$  }
\end{prooftree} 

sarebbe ancora vero il teorema di safety?
\subparagraph*{Svolgimento}
Si, lo si pu\`o dimostrare considerando che $APP^{\prime}$ non \`e altro che una specificazione di $APP$: se poniamo $X$ uguale all'insieme dei tipi di $APP$, $Y$ \`e un sottoinsieme di $X$ corrispondente ai possibili tipi di $APP^{\prime}$.\\
In sostanza la regola restringere l'insieme dei programmi realizzabili con il precedente type system

Di conseguenza il teorema di safety vale anche per $APP^{\prime}$ poich\`e la dimostrazioni continuano a valere.

\subsubsection*{Esercizio 3.15}
Se al posto delle regole $(APP)\:$ e $(FUN)\:$ si definissero le seguenti regole:
\begin{prooftree} 
	\AxiomC{$\Gamma\:\vdash$ $M\: :  $T$\:\rightarrow\:$T}
	\AxiomC{$\Gamma\:\vdash$ $N\ :$T}
	
	\LeftLabel{\textsc{(APP')}}
	\BinaryInfC{$\Gamma\:\vdash$ $M$ $N$ : $T$  }
\end{prooftree} 

e


\begin{prooftree} 
	\AxiomC{$\Gamma\:$, x : $T1\vdash$ $M\: : $T}
	\AxiomC{}
	\LeftLabel{\textsc{(FUN')}}
	\BinaryInfC{$\Gamma\:\vdash$ $fn$ x : $T1.M:T\rightarrow\:$T  } 
\end{prooftree} 

sarebbe ancora vero il teorema di safety?

\subparagraph*{Svolgimento}

La risposta e' no:\\
In questo caso la $FUN_1$ non fa alcun controllo sul tipo dell'argomento e quindi esiste la possibilit\`a di arrivare ad un passo della derivazione che produce un termine STUCK, invalidando il teorema stesso. Questo perch\`e la nuova regola risulta essere pi\`u stringente rispetto a prima visto che ammette solo tipi uguali al tipo di ritorno senza controllare effettivamente che la signatura dei paremtri sia conforme a quanto richiesto.\\
Supponiamo di avere il termine:\\
$M$ $=$ $(fn$ $x$:$Bool.if$ $x$ $then$ $1$ $else$ $0)$ $(1)$ 

e creiamo il suo albero di derivazione:




\begin{prooftree} 
	\AxiomC{$ \checkmark $}
	\UnaryInfC{$x$:$Bool$ $\in$  $\Gamma$ }
	\RightLabel{\textsc{($Var$)}}
	\UnaryInfC{$\Gamma$ $\vdash$ $x$:$Bool$ }
	\AxiomC{$ \checkmark $}
	\LeftLabel{\textsc{($Nat$)}}
	\UnaryInfC{$\Gamma$ $\vdash$ $1$:$Nat$}
	\AxiomC{$ \checkmark $}
	\LeftLabel{\textsc{($Nat$)}}
	\UnaryInfC{$\Gamma$ $\vdash$ $0$:$Nat$}
	\LeftLabel{\textsc{($IfThenElse$)}}
	\TrinaryInfC{$x$:$Bool$ $\vdash$ $if$ $x$ $then$ $1$ $else$ $0$ :$Nat$}
	\LeftLabel{\textsc{($FUN_1$)}}
	\UnaryInfC{$\empty \vdash fn$ $x$:$Bool.if$ $x$ $then$ $1$ $else$ $0$ :$Nat \rightarrow$ :$Nat$}
	\AxiomC{$ \checkmark $}
	\LeftLabel{\textsc{($Nat$)}}
	\UnaryInfC{$\empty \vdash 1$:$Nat$}
	\LeftLabel{\textsc{($APP_1$)}}
	\BinaryInfC{$\empty \vdash (fn$ $x$:$Bool.if$ $x$ $then$ $1$ $else$ $0)$ $(1)$  } 
\end{prooftree}


In questo caso possiamo osservare che il termine $M$ risulta essere ben-tipato secondo le nuove regole di tipo.\\
Tuttavia, \`e evidente che l'invocazione della funzione in questione porta $M$ ad un termine stuck, questo perch\`e il sottotermine $N$ $=$ $if$ $x$ $then$ $1$ $else$ $0$ si aspetta un $Bool$, infatti le regole della semantica operazionale che gestiscono l'$if$, permettono di avanzare nella computazione.\\
Nel caso in cui la guardia non sia un valore, la computazione procede con la valutazione della guardia, altrimenti procede il ramo $TRUE$ o con il ramo $FALSE$.\\
In questo caso, poich\`e il valore passato alla funzione \`e un intero, e quindi gi\`a un valore le uniche regole applicabili sarebbero $(IF-TRUE)$ e $(IF-FALSE)$ ma poich\`e la guardia \'e 1 non sono applicabili e quindi si otterrebbe un termine STUCK.


\subsubsection*{Esercizio 3.16}
Se si aggiungessero al sistema di tipi i seguenti due assiomi

\begin{prooftree} 
	\AxiomC{}
	\AxiomC{}
	\LeftLabel{\textsc{(TRUE')}}
	\BinaryInfC{$\Gamma\:\vdash$ true : $Nat$  } 
\end{prooftree}

e

\begin{prooftree} 
	\AxiomC{}
	\AxiomC{}
	\LeftLabel{\textsc{(FALSE')}}
	\BinaryInfC{$\Gamma\:\vdash$ false : $Nat$  } 
\end{prooftree}

sarebbe ancora vero il teorema di safety?
\subparagraph*{Svolgimento}
La risposta e' no:\\
Data la definizione di Teorema di Safety, e' evidente che viene meno la stessa ipotesi ovvero che M sia un valore o un che non evolva in un termine STUCK.

Con queste regole di tipo, infatti, sarebbe permessa anche l'operazione SUM che genererebbe un termine STUCK, invalidando il teorema stesso.\
Supponiamo di avere il termine: $M$ $=$ $true+4$


\begin{prooftree} 
	\AxiomC{$ \checkmark $}
	\LeftLabel{\textsc{($Nat$)}}
	\UnaryInfC{$\empty \vdash 4$:$Nat$}
	\AxiomC{$ \checkmark $}
	\LeftLabel{\textsc{($Nat$)}}
	\UnaryInfC{$\empty \vdash true$:$Nat$}
	\LeftLabel{\textsc{($SUM$)}}
	\BinaryInfC{$\empty \vdash$ $true$ $+$ $4$  } 
\end{prooftree}

Anche in questo caso, come nel precendete, notiamo che le regole di tipo permettono la derivazione dell'albero.\\
Tuttavia, se applicassimo le regole della semantica operazionale non potremo avanzare perch\`e entrambi gli addendi sono gi\'a dei valori e quindi non posso applicare $SUM-LEFT$ o $SUM-RIGHT$ ma non posso neppure applicare $SUM$ perch\'e genero un termine STUCK in quanto $true$ non pu\'o essere sommato a $4$.

\subsubsection*{Esercizio 3.17}

Dimostrare il seguente fatto: se $\Gamma\:\vdash$ $M$ : $T$ \`e derivabile, allora  \textit{fv($M$)}  $\subseteq$ \textit{Dom($\Gamma$)}
\subparagraph*{Svolgimento}

Si procede con la dimostrazione per induzione:

\begin{itemize}


\item Caso Base 
		\begin{itemize}[label=$\star$]

		\item \textbf{True}  $\Gamma \vdash{} M : T$ \'e $\Gamma \vdash{} \true{} : \Bool$.
  		Il teorema \'e vacuamente vero per questa derivazione siccome $\not \exists{M'}. M \to{} M'$.\\
  		Il tutto si pu\`o riassumere con $fv(M)$ $=$ $fv(true)$ $=$ $\emptyset$ $\subseteq$ $\emph{Dom(}\Gamma\emph{)}$

  
 		\item \textbf{False} $\Gamma \vdash{} M : T$ \'e $\Gamma \vdash{} \false{} : \Bool$. Analogo a True.

		\item \textbf{Nat} $\Gamma \vdash{} M : T$ \'e $\Gamma \vdash{} n : \Nat$. Analogo a True.

		\item \textbf{Var}  Se $\Gamma$ $\vdash$ $x$ : $T$ , allora per la regola $(Var)$ possiamo asserire che $x$ : $T$ $\in$ $\Gamma$ e di conseguenza abbiamo che  $fv(M)$ $=$ $fv(x)$ $=$ $\{x\}$ $\subseteq$ $\emph{Dom(}\Gamma\emph{)}$

  
		\end{itemize}
		In questo passo il fatto e' un assioma, questo perche' non esistono variabili libere quindi $fv(M)$ e' $\emptyset$ e il vuoto e un sottoinsieme del $Dom(\Gamma)$

\item Caso Induttivo
		
		\begin{itemize}[label=$\star$]

		\item \textbf{SUM}  in questo caso $M$ $\equiv$ $A$ $+$ $B$ e quindi per definizione le variabili libere di $M$ sono $f$ $v(A$ $+$ $B)$ $=$ $f$ $v(A)$ $\cup$ $f$ $v(B)$.\\ Per ipotesi induttiva sappiamo che $fv(A)$ e $fv(B)$ $\subseteq$ $\emph{Dom(}\Gamma\emph{)}$.\\Poich\`e, come anticipato, $fv(M)$ $\equiv$ $fv(A)$ $\cup$ $fv(B)$ possiamo dedurre che $fv(M)$ $\subseteq$ $\emph{Dom(}\Gamma\emph{)}$.
  		
  
 		\item \textbf{MINUS} in questo caso $M$ $\equiv$ $A$ $-$ $B$ e la sua dimostrazione risulta essere analoga a quella di SUM

		\item \textbf{IfThenElse} in questo caso $M$ $\equiv$ $if$ $C$ $then$ $A$ $else$ $B$ e quindi per definizione, le variabili libere di $M$ sono $fv(M)$ $=$ $fv(A)$ $\cup$ $fv(B)$ $\cup$ $fv(C)$.\\ Per ipotesi induttiva sappiamo che $fv(A)$, $fv(B)$ e $fv(C)$ $\subseteq$ $\emph{Dom(}\Gamma\emph{)}$.\\Poich\`e, come anticipato, $fv(M)$ $\equiv$ $fv(A)$ $\cup$ $fv(B)$ $\cup$ $fv(C)$ possiamo dedurre che $fv(M)$ $\subseteq$ $\emph{Dom(}\Gamma\emph{)}$.

		\item \textbf{FUN}  in questo caso $M$ $\equiv$ $fn$ $x.N$ e quindi per definizione le variabili libere di $M$ sono 
		$fv(M)$ $=$ $fv(N)\textbackslash \{x\}$ ($fv(M)$ $\subseteq$ $fv(N)$).\\ Per ipotesi induttiva sappiamo che $fv(N)$ $\subseteq$ $\emph{Dom(}\Gamma\emph{)}$.\\Poich\`e, come anticipato, $fv(M)$ $\equiv$ $fv(N)\textbackslash \{x\}$ possiamo dedurre che $fv(M)$ $\subseteq$ $\emph{Dom(}\Gamma\emph{)}$. 
		
		\item \textbf{APP} in questo caso $M$ $\equiv$ $A$ $B$ e quindi per definizione, le variabili libere di $M$ sono $fv(M)$ $=$ $fv(A)$ $\cup$ $fv(B)$.\\ Per ipotesi induttiva sappiamo che $fv(A)$ e $fv(B)$ $\subseteq$ $\emph{Dom(}\Gamma\emph{)}$.\\Poich\`e, come anticipato, $fv(M)$ $\equiv$ $fv(A)$ $\cup$ $fv(B)$ possiamo dedurre che $fv(M)$ $\subseteq$ $\emph{Dom(}\Gamma\emph{)}$. 
  
		\end{itemize}

\end{itemize}

\subsubsection*{Esercizio 3.18}
Ricostruire il tipo dei seguenti termini:
\begin{itemize}
\item $fn$ $x$:$T1.fn$ $y$:$T2.if$ $y$ then $x$ else true
\item $fn$ $x:Nat:\rightarrow\:$Bool.x
\item $fn$ $f:T.fn$ $x:T'.f$( if true  then  $x$  else  $f$ $x$)
\item $fn$ $f:T1.fn$ $g:T2.$ if ($f$ ($g$ true)) then $f$ ($fn$ $x:$T3.true) else $f$($fn$ $x:T4.x$)
\end{itemize}

\subparagraph*{Svolgimento}
\begin{enumerate}[label=\textbf{\alph*}), leftmargin=*]


\item $fn$ $x$:$T_1.fn$ $y$:$T_2.if$ $y$ then $x$ else true  \\
		\scalebox{.75}{
		\parbox{1cm}{
		\begin{prooftree}
			\AxiomC{y:Bool $\in$ $\Gamma$}
			\LeftLabel{\textsc{(VAR)}} 
			\LeftLabel{\textsc{(VAR)}}
			\UnaryInfC{$\Gamma$,$x$:$T_1$,$y$:$T_2$ $\vdash$ $y$ : $Bool$}
			\LeftLabel{\textsc{(VAR)}}
			\AxiomC{$x$:$T_3$ $\in$ $\Gamma$}
			\LeftLabel{\textsc{(VAR)}}
			\UnaryInfC{$\Gamma$,$x$:$T_1$,$y$:$T_2$ $\vdash$ $x$ : $T_3$}
			\LeftLabel{\textsc{(VAR)}}
			\AxiomC{$\checkmark$}
			\LeftLabel{\textsc{(TRUE)}}
			\UnaryInfC{$\Gamma$,$x$:$T_1$,$y$:$T_2$ $\vdash$ $true$ : $Bool$}
			\LeftLabel{\textsc{(IF-THEN-ELSE)}}
			\TrinaryInfC{$x$:$T_1$,$y$:$T_2$ $\vdash$ $if$ $y$ $then$ $x$ $else$ $true$ : $T_3$}
			\LeftLabel{\textsc{(FUN)}}
			\UnaryInfC{$x$:$T_1$ $\vdash$  $fn$ $y$:$T_2.if$ $y$ $then$ $x$ $else$ $true$ : $T_2 \rightarrow T_3$}
			\LeftLabel{\textsc{(FUN)}}
			\UnaryInfC{$\emptyset\:\vdash$  $fn$ $x$:$T_1.fn$ $y$:$T_2.if$ $y$ $then$ $x$ $else$ $true$ : $T_1 \rightarrow T_2 \rightarrow T_3$}
		\end{prooftree}}
		}\vspace{1cm}
		
		Il risultato finale \`e dato dal fatto che $T_1$ = $T_2$ = $T_3$ = $Bool$  e quindi il tipo finale risulta: $Bool \rightarrow Bool \rightarrow Bool$
		\vspace{1.5cm}
		
		\item $fn$ $x:Nat:\rightarrow$ $Bool.x$  \\
		\scalebox{.75}{
		\parbox{1cm}{
		\begin{prooftree}
		
			\AxiomC{$x:T_2$ $\in$ $\Gamma,$ $x$:$Nat:\rightarrow Bool$}
			\LeftLabel{\textsc{(VAR)}}
			\UnaryInfC{$\Gamma,$ $x$:$Nat:\rightarrow Bool$ $\vdash$  $x$:$T_2$}
			\LeftLabel{\textsc{(FUN)}}
			\UnaryInfC{$\emptyset\:\vdash$  $fn$ $x:Nat:\rightarrow$ $Bool.x$ : $T_1 \rightarrow T_2$}
		\end{prooftree}}
		}\vspace{1cm}
		
		Il risultato finale \`e dato dal fatto che $T_2$ = $Nat \rightarrow Bool$  e quindi il tipo finale risulta $(Nat \rightarrow Bool)  \rightarrow (Nat \rightarrow Bool)$
		\vspace{1.5cm}
		
		
		\item $fn$ $f:T.fn$ $x:T'.f(if$ $true$  $then$  $x$  $else$  $f$ $x)$ \\ 
		\scalebox{.75}{
		\parbox{1cm}{ \vspace{1.5cm}
		\begin{prooftree}
			
			%\AxiomC{$f$:$T$,$x$:$T_1$ $\vdash$ $f$:$T_3 \rightarrow T_2$}
			
			%%else
			
			%\AxiomC{$f$:$T$,$x$:$T_1$ $\vdash$ $f:T_4 \rightarrow T$}
			%\AxiomC{$T_1 = T_4$}
			%\LeftLabel{\textsc{(VAR)}}
			%\UnaryInfC{$f$:$T$,$x$:$T_1$ $\vdash$ $x:T_4$}
			%\LeftLabel{\textsc{(APP)}}
			%\BinaryInfC{$f$:$T$,$x$:$T_1$ $\vdash$ $f x$ : $T_3$}
			
			%%			
					
			%%then
			
			%\AxiomC{$T_1 = T_3$}
			%\LeftLabel{\textsc{ }}
			%\UnaryInfC{$f$:$T$,$x$:$T_1$ $\vdash$ $x$ : $T_3$}
						
			%%			
			%% if
			%\AxiomC{$\checkmark$}
			%\LeftLabel{\textsc{(TRUE)}}
			%\UnaryInfC{$f$:$T$,$x$:$T_1$ $\vdash$ $true$ : $Bool$}
			%%
			%\LeftLabel{\textsc{(IF-THEN-ELSE)}}
			%\TrinaryInfC{$f$:$T$,$x$:$T_1$ $\vdash$ $ if true$ then $x$ else $f$ $x$ : $T_3$}
			
			\AxiomC{A}
			\UnaryInfC{$f$:$T$,$x$:$T^{\prime}$ $\vdash$ $f$:$^{\prime\prime\prime\prime} \rightarrow T^{\prime\prime}$}
			\AxiomC{B}
			\UnaryInfC{$f$:$T$,$x$:$T^{\prime}$ $\vdash$ $ if true$ then $x$ else $f$ $x$ : $T_3$}
			\LeftLabel{\textsc{(APP)}}
			\BinaryInfC{$f$:$T$,$x$:$T^{\prime}$ $\vdash$ $f$ $(if$ $true$ $then$ $x$ $else$ $f$ $x$)}
			\LeftLabel{\textsc{(FUN)}}
			\UnaryInfC{$f$:$T$ $\vdash$  $fn$ $x$:$T^{\prime}.f(if$ $true$  $then$  $x$  $else$  $f$ $x)$:$T^{\prime}\rightarrow T^{\prime\prime} $}
			\LeftLabel{\textsc{(FUN)}}
			\UnaryInfC{$\emptyset\:\vdash$  $fn$ $f:T.fn$ $x:T^{\prime}.f(if$ $true$  $then$  $x$  $else$  $f$ $x)$ : $T \rightarrow T^{\prime}\rightarrow T^{\prime\prime} $}
		\end{prooftree}}
		}\vspace{1cm}
		
		
		
		\begin{itemize}[label=$\star$]
		
		\item 	Albero A \\
		
		\scalebox{.75}{
		\parbox{1cm}{\vspace{1cm}
		\begin{prooftree}
			\AxiomC{ $f$:$T^{\prime\prime\prime\prime} \rightarrow T^{\prime\prime}$ $\vdash$ $\Gamma$}
			\LeftLabel{\textsc{(VAR)}}	
			\UnaryInfC{$f$:$T$,$x$:$T^{\prime}$ $\vdash$ $f$:$T^{\prime\prime\prime\prime} \rightarrow T^{\prime\prime}$}
		
		\end{prooftree}}	
		}	\vspace{1cm}
		
		\item Albero B\\
		
		\scalebox{.75}{
		\parbox{1cm}{\vspace{1cm}
		\begin{prooftree}
		
		
			%%			
			%% if
			\AxiomC{$\checkmark$}
			\LeftLabel{\textsc{(TRUE)}}
			\UnaryInfC{$f$:$T$,$x$:$T^{\prime}$ $\vdash$ $true$ : $Bool$}
			
			
			%%			
					
			%%then
			
			\AxiomC{$x$ : $T^{\prime\prime}$ $\in$ $\Gamma$}
			\LeftLabel{\textsc{(VAR) }}
			\UnaryInfC{$f$:$T$,$x$:$T^{\prime}$ $\vdash$ $x$ : $T^{\prime\prime}$}
						
			%%else
			
			
			\AxiomC{C}
			\LeftLabel{\textsc{(APP)}}
			\UnaryInfC{$f$:$T$,$x$:$T^{\prime}$ $\vdash$ $f x$ : $T^{\prime\prime}$}
			
			%%
			\LeftLabel{\textsc{(IF-THEN-ELSE)}}
			\TrinaryInfC{$f$:$T$,$x$:$T^{\prime}$ $\vdash$ $if$ $true$ $then$ $x$ $else$ $f$ $x$ : $T^{\prime\prime}$}
		
		\end{prooftree}}	
		}	\vspace{1cm}
	
		
		\item Albero C (ramo else dell'if)\\
		
		\scalebox{.75}{
		\parbox{1cm}{\vspace{1cm}
		\begin{prooftree}
		
			\AxiomC{$f:T^{\prime\prime\prime}$ $\in$ $\Gamma$}
			\LeftLabel{\textsc{(VAR)}}
			\UnaryInfC{$f$:$T$,$x$:$T^{\prime}$ $\vdash$ $f:T^{\prime\prime\prime} \rightarrow T^{\prime\prime}$}
			\AxiomC{$x:T^{\prime\prime\prime}$ $\in$ $\Gamma$}
			\LeftLabel{\textsc{(VAR)}}
			\UnaryInfC{$f$:$T$,$x$:$T^{\prime}$ $\vdash$ $x:T^{\prime\prime\prime}$}
			\BinaryInfC{$f$:$T$,$x$:$T^{\prime}$ $\vdash$ $f x$ : $T^{\prime\prime}$}		
		
		\end{prooftree}}	
		}	\vspace{1cm}	
		
		
		\end{itemize}
		
		il risultato finale \`e supportato dalle seguenti sostituzioni:
		\begin{itemize}
		
		\item $T \rightarrow T^{\prime} \rightarrow T^{\prime\prime}$
		\item $T^{\prime} \rightarrow T^{\prime\prime}$
		\item $T^{\prime\prime\prime} \rightarrow T^{\prime\prime}$
		\item $T^{\prime}$ =  $T^{\prime\prime\prime}$
		\item $T$ = $T^{\prime\prime\prime\prime}$
		\item $T^{\prime\prime\prime\prime}$ = $T^{\prime\prime\prime}$
		\item $T$ =  $T^{\prime\prime\prime}$
		\item $T^{\prime\prime}$  =  $T^{\prime\prime\prime}$
			
		
		\end{itemize}
				
		Il risultato finale e' dato dal fatto che $T = T_1 = T_" = T_3 = T_4$ quindi $(T \rightarrow T) \rightarrow (T \rightarrow T) $
		\vspace{1.5cm}
		
	
\item $fn$ $f:T1.fn$ $g:T2.if$ $(f$ $(g$ $true))$ $then$ $f$ ($fn$ $x:T3.true)$ $else$ $f(fn$ $x:T4.x)$  \\
		
		\scalebox{.75}{
		\parbox{1cm}{\vspace{1.5cm}
		\begin{prooftree}
			
				
			%guardia if
			\AxiomC{A}
			\UnaryInfC{$f$:$T_1$,$g$:$T_2$ $\vdash$ $f$ $(g$ $true))$ : $Bool$}	
			
			%then
			\AxiomC{B}
			\UnaryInfC{$f$:$T_1$,$g$:$T_2$ $\vdash$ $f$ $(fn$ $x$:$T_3.true)$:$W$}
			
			%else
			\AxiomC{C}
			\UnaryInfC{$f$:$T_1$,$g$:$T_2$ $\vdash$ $f(fn$ $x:T_4.x)$:$W$}
			
						
			
			\LeftLabel{\textsc{(IF-THEN-ELSE)}}
			\TrinaryInfC{$f$:$T_1$,$g$:$T_2$ $\vdash$ $if$ $(f$ $(g$ $true))$ $then$ $f$ ($fn$ $x:T_3.true)$ $else$ $f(fn$ $x:T_4.x)$:$W$}
			\LeftLabel{\textsc{(FUN)}}
			\UnaryInfC{$f$:$T_1$ $\vdash$ $fn$ $g:T_2.if$ $(f$ $(g$ $true))$ $then$ $f$ ($fn$ $x:T_3.true)$ $else$ $f(fn$ $x:T_4.x)$:$T_2 \rightarrow W$}
			\LeftLabel{\textsc{(FUN)}}
			\UnaryInfC{$\emptyset\:\vdash$  $fn$ $f:T_1.fn$ $g:T_2.if$ $(f$ $(g$ $true))$ $then$ $f$ ($fn$ $x:T_3.true)$ $else$ $f(fn$ $x:T_4.x)$:$T_1 \rightarrow T_2 \rightarrow W$}
		\end{prooftree}}
		}	
		
		\vspace{1cm}
		Dove gli alberi A, B e C sono:
		

		\begin{itemize}[label=$\star$]
		
		\item 	Albero A (derivazione guardia del'if)\\
		
		\scalebox{.75}{
		\parbox{1cm}{\vspace{1cm}
		\begin{prooftree}
		
			\AxiomC{$f$:$ U \rightarrow Bool $ $\in$ $f$:$T_1$,$g$:$T_2$}
			\LeftLabel{\textsc{(VAR)}}
			\UnaryInfC{$f$:$T_1$,$g$:$T_2$ $\vdash$ $f$:$ U \rightarrow Bool $}
			
			
			\AxiomC{$g$:$ Bool \rightarrow U $ $\in$ $f$:$T_1$,$g$:$T_2$}
			\LeftLabel{\textsc{(VAR)}}
			\UnaryInfC{$f$:$T_1$,$g$:$T_2$ $\vdash$ $g$:$ Bool \rightarrow U $}
			\AxiomC{$\checkmark$}
			\LeftLabel{\textsc{(TRUE)}}
			\UnaryInfC{$f$:$T_1$,$g$:$T_2$ $\vdash$ $true$:$ Bool $}
			\LeftLabel{\textsc{(APP)}}
			\BinaryInfC{$f$:$T_1$,$g$:$T_2$ $\vdash$ $g$ $true$:$U$}
			
			\LeftLabel{\textsc{(APP)}}
			\BinaryInfC{$f$:$T_1$,$g$:$T_2$ $\vdash$ $f$ $(g$ $true))$ : $Bool$}
		\end{prooftree}}	
		}	\vspace{1cm}
		
		
		\item 	Albero B (derivazione ramo then)\\
		
		\scalebox{.75}{
		\parbox{1cm}{\vspace{1cm}
		\begin{prooftree}
		
			\AxiomC{$f: Z \rightarrow W$ $\in$ $f$:$T_1$,$g$:$T_2$}
			\LeftLabel{\textsc{(VAR)}}	
			\UnaryInfC{$f$:$T_1$,$g$:$T_2$ $\vdash$ $f$ $Z \rightarrow W $}
			
			\AxiomC{$\checkmark$}
			\LeftLabel{\textsc{(TRUE)}}
			\UnaryInfC{$f$:$T_1$,$g$:$T_2$,$x$:$T_3$ $\vdash$ $true$ : $Bool$}
			\LeftLabel{\textsc{(FUN)}}
			\UnaryInfC{$f$:$T_1$,$g$:$T_2$ $\vdash$ $fn$ $x$:$T_3.true$ : $T_3 \rightarrow Bool$}
			
			\LeftLabel{\textsc{(APP)}}
			\BinaryInfC{$f$:$T_1$,$g$:$T_2$ $\vdash$ $f$ $(fn$ $x$:$T_3.true)$ :$W$}
		\end{prooftree}}	
		}\vspace{1cm}			
		
		
		
		\item 	Albero C (derivazione ramo else)\\
		
		\scalebox{.75}{
		\parbox{1cm}{\vspace{1cm}
		\begin{prooftree}
		
			\AxiomC{$f: Z_1 \rightarrow W$ $\in$ $f$:$T_1$,$g$:$T_2$}
			\LeftLabel{\textsc{(VAR)}}
			\UnaryInfC{$f$:$T_1$,$g$:$T_2$ $\vdash$ $f$ $Z_1 \rightarrow W $}		
		
			\AxiomC{$x$:$T_4$ $\in$ $\Gamma$}
			\LeftLabel{\textsc{(VAR)}}
			\UnaryInfC{$f$:$T_1$,$g$:$T_2$,$x$:$T_4$ $\vdash$ $x$ : $T_4$}
			\LeftLabel{\textsc{(FUN)}}
			\UnaryInfC{$f$:$T_1$,$g$:$T_2$ $\vdash$ $fn$ $x$:$T_4.x$ : $T_4 \rightarrow T_4$}
			
			
			\LeftLabel{\textsc{(APP)}}
			\BinaryInfC{$f$:$T_1$,$g$:$T_2$ $\vdash$ $f(fn$ $x:T_4.x)$:$W$}
		\end{prooftree}}	
		}\vspace{1cm}					
		
		
		\end{itemize}				
		
		Il risultato finale \`e supportato dalle seguenti sostituzioni:
		\begin{itemize}
		
		\item $T_1$ = $U \rightarrow Bool$
		\item $T_1$ = $Z \rightarrow Bool$
		\item $T_1$ = $Z_1 \rightarrow Bool$
		\item $T_2$ = $Bool \rightarrow U$
		\item $W$ = $Bool$
		\item $U$ = $Bool \rightarrow Bool$
		\item $T_3$ = $T_4$ = $Bool$
		
		\item quindi $T_2$ = $Bool \rightarrow (Bool \rightarrow Boool)$ ( per sostituzione infatti $f$ $(fn$ $x$:$T_3.true)$ \`e $U \rightarrow W$ con $W$ = $Bool$ e $U$ = $(Bool \rightarrow Bool))$
		\item quindi $T_1$ =  $(Bool \rightarrow Bool) \rightarrow Boool$
		\item quindi da $T_1 \rightarrow T_2 \rightarrow W$ otteniamo $(Bool \rightarrow Bool) \rightarrow Boool$ $\rightarrow$ $(Bool \rightarrow (Bool \rightarrow Boool))$ $\rightarrow$ $Bool$
		
		
		
		\end{itemize}
		e quindi il tipo del termine \`e quindi:\\  $(Bool \rightarrow Bool) \rightarrow Boool$ $\rightarrow$ $(Bool \rightarrow (Bool \rightarrow Boool))$ $\rightarrow$ $Bool$
		

\end{enumerate}



\newpage
\bibliographystyle{plain}
\bibliography{biblist}

\end{document} 
